\documentclass[final]{beamer}

\usepackage[scale=1.24, size=a1, orientation=portrait]{beamerposter}
\usepackage[utf8]{inputenc}
\usepackage[T1]{fontenc}
\usepackage[french]{babel}
\usepackage{amsmath,amssymb}
\usepackage{booktabs}
\usepackage{tikz}
\usepackage{xcolor}
\usepackage{graphicx}
\usepackage{multicol}
\usepackage{enumitem}

% ─── Couleurs ────────────────────────────────────────────────────────────────
\definecolor{mainblue}{RGB}{0,84,166}
\definecolor{lightblue}{RGB}{210,228,255}
\definecolor{darkgray}{RGB}{50,50,50}
\definecolor{accentorange}{RGB}{230,100,20}
\definecolor{greentick}{RGB}{30,160,80}
\definecolor{blockbg}{RGB}{240,245,255}

% ─── Thème Beamer ────────────────────────────────────────────────────────────
\usetheme{default}
\setbeamercolor{block title}{fg=white,bg=mainblue}
\setbeamercolor{block body}{fg=darkgray,bg=blockbg}
\setbeamercolor{block title alerted}{fg=white,bg=accentorange}
\setbeamercolor{block body alerted}{fg=darkgray,bg=blockbg}
\setbeamerfont{block title}{size=\large,series=\bfseries}
\setbeamerfont{block body}{size=\normalsize}
\setbeamertemplate{navigation symbols}{}

% ─── En-tête personnalisé ────────────────────────────────────────────────────
\setbeamertemplate{headline}{
  \leavevmode
  \begin{beamercolorbox}[wd=\paperwidth]{headline}
    \begin{tikzpicture}[remember picture]
      \fill[mainblue] (0,0) rectangle (\paperwidth, 7cm);
      \node[anchor=center, text=white] at (0.5\paperwidth, 5cm)
        {\fontsize{52}{60}\bfseries Détection de Deepfakes};
      \node[anchor=center, text=white!90] at (0.5\paperwidth, 3.2cm)
        {\fontsize{28}{34}\selectfont Approche hybride : Vision Transformer (DINOv2) + Analyse Fréquentielle};
      \node[anchor=center, text=white!75] at (0.5\paperwidth, 1.5cm)
        {\fontsize{20}{24}\selectfont Projet M2 IA --- Université --- 2025--2026};
    \end{tikzpicture}
  \end{beamercolorbox}
}

\setbeamertemplate{footline}{
  \leavevmode
  \begin{beamercolorbox}[wd=\paperwidth, ht=1.8cm]{headline}
    \begin{tikzpicture}
      \fill[mainblue] (0,0) rectangle (\paperwidth, 1.8cm);
      \node[anchor=west, text=white] at (1cm, 0.9cm)
        {\small \textbf{Code :} github.com/anessrb/deepfakedetection};
      \node[anchor=center, text=white] at (0.5\paperwidth, 0.9cm)
        {\small Bibliothèques : PyTorch $\geq$ 2.1 \textbullet\ timm $\geq$ 0.9 \textbullet\ albumentations $\geq$ 1.3 \textbullet\ scikit-learn};
      \node[anchor=east, text=white] at (\paperwidth-1cm, 0.9cm)
        {\small 2026};
    \end{tikzpicture}
  \end{beamercolorbox}
}

% ─── Macro bloc coloré ───────────────────────────────────────────────────────
\newcommand{\cbox}[3]{%
  \begin{tikzpicture}
    \node[fill=#1, rounded corners=6pt, inner sep=8pt, text width=#2, align=center]
      {\color{white}\bfseries #3};
  \end{tikzpicture}%
}

\begin{document}
\begin{frame}[t]
\vspace{0.5cm}

% ════════════════════════════════════════════════════════════════════════════
%  LIGNE 1 : Introduction + Architecture
% ════════════════════════════════════════════════════════════════════════════
\begin{columns}[t]

%% ── Col gauche : Introduction ──────────────────────────────────────────────
\begin{column}{0.31\linewidth}

\begin{block}{Introduction}
Un \textbf{deepfake} est un contenu audiovisuel falsifié par des réseaux de neurones
(GAN, diffusion) qui peut tromper l'œil humain.\\[6pt]
\textbf{Enjeux :}
\begin{itemize}[leftmargin=1.2em, itemsep=2pt]
  \item Désinformation et manipulation politique
  \item Usurpation d'identité
  \item Atteinte à la vie privée
\end{itemize}
\vspace{6pt}
\textbf{Objectif :} concevoir un détecteur automatique robuste, calibré et interprétable,
capable de généraliser sur des datasets et des compressions variées.
\end{block}

\vspace{0.4cm}

\begin{block}{Datasets utilisés}
\centering
\small
\begin{tabular}{@{}lll@{}}
\toprule
\textbf{Dataset} & \textbf{Rôle} & \textbf{Type}\\
\midrule
FaceForensics++ & Train & 4 manipulations\\
AV-Deepfake1M  & Train & In-the-wild\\
DF40           & Train/Val & 40 méthodes\\
Celeb-DF v2    & Test & Face-swap réaliste\\
WildDeepfake   & Test & In-the-wild\\
GenFace        & Test & Diffusion\\
\bottomrule
\end{tabular}
\vspace{4pt}

Division : \textbf{70\%} train / \textbf{15\%} val / \textbf{15\%} test (seed=42)\\
Équilibrage par \textit{WeightedRandomSampler}.
\end{block}

\vspace{0.4cm}

\begin{block}{Pré-traitement}
\begin{enumerate}[leftmargin=1.4em, itemsep=3pt]
  \item Extraction de frames (1 fps, max 300)
  \item Détection de visage (OpenCV Haar Cascade)
  \item Recadrage 224×224 + marge 20\%
  \item Normalisation ImageNet
\end{enumerate}
\vspace{6pt}
\textbf{Augmentations (entraînement) :}
\begin{itemize}[leftmargin=1.2em, itemsep=2pt]
  \item RandomResizedCrop (scale 0.8--1.0)
  \item HorizontalFlip $p=0.5$
  \item ColorJitter (brightness/contrast/saturation 0.2)
  \item JPEG compression (qualité 30--100, $p=0.5$)
  \item GaussianBlur ($\sigma \in [0.1, 2.0]$, $p=0.3$)
\end{itemize}
\end{block}

\end{column}

%% ── Col centre : Architecture ──────────────────────────────────────────────
\begin{column}{0.36\linewidth}

\begin{block}{Architecture du Modèle}
\centering
\begin{tikzpicture}[
  node distance=0.55cm,
  box/.style={draw=mainblue, fill=lightblue, rounded corners=5pt,
              minimum width=5.5cm, minimum height=0.85cm,
              font=\small\bfseries, text=darkgray, align=center},
  sbox/.style={draw=accentorange, fill=accentorange!15, rounded corners=5pt,
               minimum width=2.5cm, minimum height=0.85cm,
               font=\small\bfseries, align=center},
  arr/.style={->, thick, mainblue},
  sarr/.style={->, thick, accentorange},
  ]

  % Input
  \node[box, fill=mainblue!80, text=white] (inp) {Image d'entrée $224\times224\times3$};

  % Face det
  \node[box, below=of inp] (fd) {Détection de visage (OpenCV)};
  \draw[arr] (inp) -- (fd);

  % Split
  \node[sbox, below left=0.7cm and 1.5cm of fd, text=mainblue] (sp) {Branche\\Spatiale};
  \node[sbox, below right=0.7cm and 1.5cm of fd, text=accentorange] (fr) {Branche\\Fréquentielle};
  \draw[arr] (fd.south) -- ++(0,-0.25) -| (sp.north);
  \draw[arr] (fd.south) -- ++(0,-0.25) -| (fr.north);

  % Spatial details
  \node[box, below=0.5cm of sp, minimum width=4.5cm] (dino) {DINOv2 ViT-B/14\\(4 derniers blocs)};
  \node[box, below=0.45cm of dino, minimum width=4.5cm] (cls) {CLS token\\$\mathbf{f}_s \in \mathbb{R}^{768}$};
  \draw[arr] (sp) -- (dino);
  \draw[arr] (dino) -- (cls);

  % Freq details
  \node[box, below=0.5cm of fr, minimum width=4.5cm] (fft) {FFT 2D + DCT-II\\$\to$ 2 canaux $224\times224$};
  \node[box, below=0.45cm of fft, minimum width=4.5cm] (lcnn) {LightCNN\\{[32,64,128,256]}};
  \node[box, below=0.45cm of lcnn, minimum width=4.5cm] (femb) {$\mathbf{f}_f \in \mathbb{R}^{256}$};
  \draw[arr] (fr) -- (fft);
  \draw[arr] (fft) -- (lcnn);
  \draw[arr] (lcnn) -- (femb);

  % Concat
  \node[box, below=1.9cm of fd, fill=greentick!80, text=white, minimum width=5.8cm] (cat)
    {Concaténation $[\mathbf{f}_s \| \mathbf{f}_f] \in \mathbb{R}^{1024}$};
  \draw[arr] (cls.south) -- ++(0,-0.3) -| (cat.north west);
  \draw[arr] (femb.south) -- ++(0,-0.3) -| (cat.north east);

  % MLP
  \node[box, below=0.5cm of cat] (mlp) {Fusion MLP : $1024\!\to\!512\!\to\!128\!\to\!1$};
  \draw[arr] (cat) -- (mlp);

  % Temp
  \node[box, below=0.45cm of mlp] (ts) {Temperature Scaling ($T^*$)};
  \draw[arr] (mlp) -- (ts);

  % Output
  \node[box, below=0.45cm of ts, fill=accentorange, text=white] (out)
    {$P(\text{fake}) \in [0,1]$};
  \draw[arr] (ts) -- (out);

\end{tikzpicture}
\end{block}

\vspace{0.4cm}

\begin{alertblock}{Branche Fréquentielle — Détail}
Les deepfakes laissent des \textbf{artefacts invisibles à l'œil} mais détectables en fréquence.

\vspace{6pt}
\textbf{FFT 2D} (spectre log-magnitude centré) :
$$S_{\text{FFT}}(u,v) = \log\!\bigl(1 + |\mathcal{F}\{I_{\text{gray}}\}(u,v)|\bigr)$$

\textbf{DCT-II} (via extension symétrique + FFT) :
$$S_{\text{DCT}}(k) = \sum_{n=0}^{N-1} x_n \cos\!\Bigl(\tfrac{\pi k}{2N}(2n+1)\Bigr)$$

Les deux cartes ($224\!\times\!224$) sont empilées en une entrée \textbf{2 canaux} pour le LightCNN.
\end{alertblock}

\end{column}

%% ── Col droite : Entraînement ──────────────────────────────────────────────
\begin{column}{0.31\linewidth}

\begin{block}{Entraînement}
\textbf{Optimiseur :} AdamW avec taux d'apprentissage différentiels :
\vspace{4pt}
\begin{center}
\begin{tabular}{@{}lc@{}}
\toprule
\textbf{Groupe} & \textbf{LR}\\
\midrule
Backbone (ViT blocs) & $10^{-5}$\\
Tête (LightCNN + MLP) & $10^{-4}$\\
\bottomrule
\end{tabular}
\end{center}

\vspace{6pt}
\textbf{Scheduler :} LinearWarmup (2 epochs) + CosineAnnealing ($\eta_{\min}=10^{-6}$)\\[4pt]
\textbf{Précision mixte :} AMP (autocast CUDA)\\[4pt]
\textbf{Gradient clipping :} $\|\nabla\|_2 \leq 1.0$\\[4pt]
\textbf{Early stopping :} patience = 5 epochs (sur AUC val)

\vspace{8pt}
\textbf{Fonction de perte combinée :}
$$\mathcal{L} = \mathcal{L}_{\text{Focal}} + \lambda_{\text{ECE}} \cdot \mathcal{L}_{\text{ECE}}$$

\textit{Focal Loss} (Lin et al., 2017), $\gamma=2$, $\alpha=0.25$ :
$$\mathcal{L}_{\text{Focal}} = -\alpha_t(1-p_t)^\gamma \log(p_t)$$

\textit{ECE Loss} différentiable (15 bins), $\lambda=0.1$ :
$$\mathcal{L}_{\text{ECE}} = \sum_{m=1}^{M} \frac{|B_m|}{n}\bigl|\overline{\text{conf}}(B_m) - \overline{\text{acc}}(B_m)\bigr|$$
\end{block}

\vspace{0.4cm}

\begin{block}{Calibration — Temperature Scaling}
Après l'entraînement, on apprend un scalaire $T^*$ par minimisation de la NLL
sur la validation (LBFGS) :
$$p_{\text{cal}} = \sigma\!\Bigl(\frac{z}{T^*}\Bigr), \quad T^* = \arg\min_T \mathcal{L}_{\text{NLL}}\!\Bigl(\tfrac{z}{T}, y\Bigr)$$
$T>1$ réduit la confiance ; $T<1$ l'augmente. La courbe de fiabilité
doit se rapprocher de la diagonale après calibration.
\end{block}

\vspace{0.4cm}

\begin{block}{Métriques d'évaluation}
\begin{center}
\small
\begin{tabular}{@{}ll@{}}
\toprule
\textbf{Métrique} & \textbf{Description}\\
\midrule
AUC-ROC & Aire sous la courbe ROC\\
ECE & Erreur de calibration attendue\\
Accuracy & Taux de bonne classification\\
AP & Aire sous la courbe Précision-Rappel\\
TPR@95 & Seuil pour 95\% de vrai-positifs\\
\bottomrule
\end{tabular}
\end{center}

\vspace{6pt}
\textbf{Évaluation de robustesse :}
\begin{itemize}[leftmargin=1.2em, itemsep=2pt]
  \item JPEG : qualités [30, 50, 70, 90, 100]
  \item Flou Gaussien : $\sigma \in \{0, 1, 2, 3, 5\}$
  \item Redimensionnement : échelles [0.25, 0.5, 0.75, 1.0]
\end{itemize}
\end{block}

\end{column}
\end{columns}

\vspace{0.6cm}

% ════════════════════════════════════════════════════════════════════════════
%  LIGNE 2 : Résultats + Explicabilité + Perspectives
% ════════════════════════════════════════════════════════════════════════════
\begin{columns}[t]

%% ── Résultats attendus ─────────────────────────────────────────────────────
\begin{column}{0.31\linewidth}
\begin{block}{Résultats attendus (cross-dataset)}
\centering
\small
\begin{tabular}{@{}lccc@{}}
\toprule
\textbf{Dataset test} & \textbf{AUC} & \textbf{Acc.} & \textbf{ECE}\\
\midrule
FaceForensics++ & $\sim$0.98 & $\sim$0.95 & $<$0.03\\
Celeb-DF v2     & $\sim$0.92 & $\sim$0.88 & $<$0.05\\
DF40            & $\sim$0.89 & $\sim$0.85 & $<$0.06\\
WildDeepfake    & $\sim$0.87 & $\sim$0.83 & $<$0.07\\
\bottomrule
\end{tabular}
\vspace{6pt}

\textit{Valeurs indicatives — à compléter après entraînement.}
\vspace{6pt}

\textbf{Hyperparamètres clés :}\\[4pt]
\begin{tabular}{@{}ll@{}}
Epochs & 30 (early stop)\\
Batch size & 32\\
Weight decay & $10^{-4}$\\
Dropout MLP & [0.3, 0.2]\\
drop\_path & 0.1\\
\end{tabular}
\end{block}
\end{column}

%% ── Explicabilité ──────────────────────────────────────────────────────────
\begin{column}{0.36\linewidth}
\begin{block}{Explicabilité — Grad-CAM sur ViT}
\textbf{Vision Transformer Grad-CAM} permet de localiser les régions faciales
qui ont le plus contribué à la décision du modèle.

\vspace{8pt}
\begin{center}
\begin{tikzpicture}
  % Fake image box
  \node[draw=accentorange, thick, rounded corners=4pt,
        minimum width=5cm, minimum height=3.5cm, fill=accentorange!10] (img)
    at (0,0) {};
  \node[font=\small\bfseries, text=accentorange] at (0, 0)
    {Image Fake\\[2pt]{\footnotesize (visage synthétisé)}};

  % Arrow
  \draw[->, ultra thick, mainblue] (2.7,0) -- (4.3,0)
    node[midway, above, font=\footnotesize, text=mainblue] {Grad-CAM};

  % Heatmap box
  \node[draw=greentick, thick, rounded corners=4pt,
        minimum width=5cm, minimum height=3.5cm, fill=greentick!10] (heat)
    at (7.1,0) {};
  \node[font=\small\bfseries, text=greentick] at (7.1, 0)
    {Carte de chaleur\\[2pt]{\footnotesize (zones suspectes)}};

  % Labels
  \node[font=\tiny, text=darkgray] at (0,-2.2) {$P(\text{fake})=0.94$};
  \node[font=\tiny, text=darkgray] at (7.1,-2.2) {Contours / texture peau};
\end{tikzpicture}
\end{center}

\vspace{4pt}
La carte est superposée à l'image originale avec $\alpha=0.5$.\\
Méthode : \textbf{Attention Rollout} + gradients sur le bloc ViT cible.
\end{block}
\end{column}

%% ── Perspectives ───────────────────────────────────────────────────────────
\begin{column}{0.31\linewidth}
\begin{block}{Perspectives}
\begin{itemize}[leftmargin=1.4em, itemsep=5pt]
  \item[\textcolor{mainblue}{\textbullet}]
    \textbf{Fine-tuning complet} du backbone DINOv2 sur données deepfakes massives
  \item[\textcolor{mainblue}{\textbullet}]
    Intégration de datasets de \textbf{diffusion} (Stable Diffusion, FLUX, Midjourney)
  \item[\textcolor{mainblue}{\textbullet}]
    \textbf{Détection vidéo temporelle} avec attention inter-frames (ViViT)
  \item[\textcolor{mainblue}{\textbullet}]
    \textbf{Déploiement API REST} avec FastAPI + conteneurisation Docker
  \item[\textcolor{mainblue}{\textbullet}]
    Étude de la \textbf{robustesse adversariale} (attaques FGSM, PGD)
  \item[\textcolor{mainblue}{\textbullet}]
    Réduction du modèle par \textbf{distillation} pour inférence mobile
\end{itemize}

\vspace{10pt}
\begin{block}{Références}
\tiny
[1] Caron et al. (2021). \textit{Emerging properties in self-supervised ViT.} ICCV.\\
[2] Lin et al. (2017). \textit{Focal Loss for Dense Object Detection.} ICCV.\\
[3] Guo et al. (2017). \textit{On calibration of modern neural networks.} ICML.\\
[4] Rossler et al. (2019). \textit{FaceForensics++.} ICCV.\\
[5] Li et al. (2020). \textit{Celeb-DF: A Large-scale Challenging Dataset.} CVPR.\\
[6] He et al. (2024). \textit{DF40: Toward Next-Generation Deepfake Detection.}
\end{block}
\end{block}
\end{column}

\end{columns}

\end{frame}
\end{document}
